\customchapter{RESUMEN}

En el registro digital del patrimonio arquitectónico se usan escáneres láser que
generan nubes de puntos con millones de elementos, pero esas nubes no traen
ninguna etiqueta: indican dónde hay materia y no qué elemento es. Hoy la
clasificación de columnas, arcos, bóvedas, cornisas y contrafuertes se hace a
mano, elemento por elemento, y no existe una herramienta que la automatice. Los
métodos actuales de segmentación semántica en tres dimensiones tampoco sirven,
porque trabajan con una lista de clases fija y obligan a que todo punto no
previsto caiga en la clase conocida más parecida. Una tesis previa desarrollada en
UTEC sobre este mismo tipo de edificio \cite{tesisutec2024} muestra el costo de
esa limitación: un modelo entrenado sobre el patio de la catedral dejó de
reconocer columnas y arcos al aplicarse a una capilla del mismo edificio, y la
única salida disponible fue etiquetar esa capilla a mano y volver a entrenar. El
objetivo de este trabajo es mejorar de forma medible la clasificación de esos
elementos no previstos.

El enfoque de descubrimiento de clases novedosas, introducido para nubes de puntos
por Riz et al. \cite{riz2023nops} y refinado por Xu et al. \cite{xu2024dasl}, sí
permite trabajar con clases que el modelo nunca vio etiquetadas. Su mejor método
actual, HyperNCD \cite{zhang2026hyperncd}, obtiene entre 11,8 y 47,5 de mIoU en
esas clases, frente a 49,8 del modelo supervisado equivalente. Al revisar el
método se encontró que construye sus prototipos a partir de descriptores
geométricos calculados a mano, como linealidad, planaridad y dispersión, y que su
propio estudio de ablación le atribuye a esa pieza solo $+4{,}2$ puntos de mIoU,
menos de la mitad de lo que aporta el razonamiento por hipergrafo. Otro trabajo
reciente, Sonata \cite{wu2025sonata}, muestra que apoyarse en ese tipo de señales
de bajo nivel perjudica la representación y que corregirlo sube el desempeño de un
descriptor congelado de 5,6 a 72,5 de mIoU. Por eso se propone un modelo híbrido
que reemplaza esos descriptores manuales por el codificador auto-supervisado de
Sonata dentro del módulo de prototipos de HyperNCD, y se mantiene sin cambios el
resto del método para que cualquier variación en la métrica se pueda atribuir a
ese reemplazo. La viabilidad de esa integración se apoya en Point Transformer V3
\cite{wu2024ptv3}, que reduce el consumo de memoria en un factor de $10{,}2$ y es
la red sobre la que Sonata está construido. La decisión de no cambiar la
arquitectura base se sostiene sobre la comparación con Stratified Transformer
\cite{lai2022stratified}, cuya ganancia por cambio de arquitectura es del mismo
orden que la del descriptor pero exige rehacer el método completo. El modelo se
evalúa sobre el levantamiento de la Catedral de Lima, de alrededor de dos
gigabytes, dividido en cinco secciones.

Se reporta primero el punto de partida, obtenido al reproducir el método original
con tres corridas independientes. Sobre esa referencia se mide la mejora del
modelo híbrido en las clases novedosas, se verifica que las clases conocidas no
se degraden y se comparan tres variantes del descriptor: el geométrico manual, el
auto-supervisado y la unión de ambos. Además de las métricas habituales se define
una medida propia de brecha cerrada, que expresa qué parte de la distancia entre
las clases conocidas y las novedosas logra recuperar el modelo, y se reporta la
métrica también en un punto intermedio para saber si la mejora viene de la
representación o del razonamiento posterior. El reporte incluye un análisis por
clase y no solo el promedio, siguiendo la evidencia de Su et al.
\cite{su2025brdl}, para quienes corregir el sesgo hacia las clases mayoritarias
aporta hasta $15{,}34$ puntos de mIoU en las clases nuevas, y declara de forma
explícita el costo en horas de GPU y memoria, siguiendo el criterio de reporte de
Robert et al. \cite{robert2023spt}.

La conclusión que se busca es si un problema demostrado en un entorno con
supervisión completa sigue siendo un problema cuando no hay etiquetas para las
clases objetivo. Cada punto de mIoU que se gana es trabajo manual de corrección
que el especialista deja de hacer, así que el valor del proyecto no está en
igualar la precisión del trabajo manual sino en automatizarlo a un nivel que
resulte útil.

\vspace{1em}
\noindent\textbf{Palabras clave:} nubes de puntos; segmentación semántica 3D;
descubrimiento de clases novedosas; aprendizaje auto-supervisado; hipergrafos;
patrimonio arquitectónico.
